% =============================================================
%  chapter2goethite_SI_body.tex
%  Shared, class-agnostic body of Chapter 2's Supplementary
%  Information (Detailed Analytical Methods through Totals).
%  Citations use \citep/\citet against the master references.bib,
%  same convention as chapter2_body.tex.
%
%  \input{} by BOTH:
%    - Chapter2/Writing/chapter2goethite_SI.tex  (standalone draft --
%      loads biblatex itself and prints its own bibliography, same
%      pattern as chapter2goethite.tex)
%    - Chapter2/chapter2_SI.tex                     (dissertation-chapter
%      version -- relies on the single \printbibliography at the end
%      of dissertation.tex; no local bibliography needed)
%  Edit this file, not the two wrappers, to change the text,
%  equations, figures, or citations.
% =============================================================

\section{Supplementary Verification: The Na\"ive/Informed Bracketing Method for the Encapsulated, Irradiated Win-06-01b Step Heat Experiment}
\label{sec:ch2-si-detailed-methods}

This supplement documents the na\"ive/informed bracketing method: an independent, more bespoke treatment of the HT-Win-06-01b 3He dataset that explicitly tracks the helium content of the ampoule step by step, rather than bounding it as in the main text's ``treatment as uncertainty'' approach (\S\ref{experimental-methods}, \S\ref{analytical-methods}). It is retained here solely as a secondary, independent verification of the main text's results and conclusions, not as the primary basis for them: the two methods' distilled 3He datasets and derived quantities agree (\S\ref{complete-degas-test}), corroborating the main text's simpler treatment. Because 4He permeation-step signals were blank level, no analogous direct check exists for 4He; 4He is therefore reduced exclusively by the bounding method throughout this chapter and does not otherwise appear in this supplement. This supplement covers only the na\"ive/informed bracketing dataset and the bespoke, step-by-step corrections applied to it (the \nameref{raw-data-and-corrections-columns} through \nameref{totals} sections below); the isothermal-permeation-fit methodology and cooling-tube numerical model that establish the ampoule's elevated permeability -- the evidentiary basis for the main text's conservative literature-permeability assumption, and a prerequisite for the corrections applied below -- are described in full in the main text (\S\ref{experimental-methods}).

\subsection{Summary}

Helium isotope measurements of evolved gasses from hydrothermal steps
heats of glass-encapsulated irradiated HT-Win-06-01b goethite of the
Flinders Ranges \citep{monteiro2017}, were acquired at Caltech on a Thermo
Fisher Helix-MC-Plus\textsuperscript{10K} mass spectrometer in peak
hopping mode from January 31\textsuperscript{st}, 2024 to March
13\textsuperscript{th}, 2024. Raw data from these measurements are found
on the first tab ``Raw Data and Corrections'' of the
``IrrWin0601bTube012624FullDataBJT.xlsx'' file. Because the sample was
encapsulated in an imperfectly helium-transmissive silica glass ampoule,
experimental steps at temperatures at or above 170°C were bracketed by
low temperature (\textless120°C ) steps where the helium in goethite at
high pH2O was effectively immobile for experimental timescales. For 3He,
the resulting low-temperature measurements were above blank for the
majority of the experiment and were entirely controlled by the process
of permeation through silica glass, and therefore were used to estimate
the amount of residual helium within the ampoule between
high-temperature experimental steps. These corrections are carried out
in the ``Raw Data and Corrections'' tab using permeability values
derived internally within the experiment from the ``Naïve Apparent
Lambdas'' tab. Diffusivities calculated using the equations of
\citet{fechtigkalbitzer1966} for 3He and 4He can be found on the
``Diffusivities'' tab. During the step heat experiment, 116 measurements
of high 3He MM gas standard were analyzed to track instrumental drift of
3He and 4He sensitivities and to quantify internal uncertainties using
reproducibility. For the final step, the glass ampoule was broken and
the remainder of the helium was degassed in a resistance furnace at
1100°C. Data from this final step can be found on the ``Totals'' tab of
``IrrWin0601bTube012624FullDataBJT1.xlsx''.

\section*{Raw Data and Corrections
Columns}\label{raw-data-and-corrections-columns}
% \section* doesn't call \sectionmark on its own (unlike numbered \section),
% so the running header would otherwise keep showing whatever numbered
% section came before it (here, the first SI section) for the entire span
% of this and the other unnumbered sections below. \markright updates the
% header directly, bypassing \sectionmark's automatic \thesection-prepending
% (which would otherwise print this section's now-frozen, stale number).
\markright{\MakeUppercase{Raw Data and Corrections Columns}}

\subsection*{Raw Data}\label{raw-data}

``Nominal (Heated) Time'' refers to the time programmed for the step
heat. These times are used for both calculating apparent diffusivities
with the equations of \citet{fechtigkalbitzer1966} and for permeation lag
corrections.

The ``Watlow Time'' column gives the last timestamp recorded in the
Watlow temperature log at the target temperature. It is always slightly
higher than the nominal time because the recorded time includes the
initial time required to heat the sample up to the programmed
temperature: typically around 30 seconds to two minutes. This column was
generated by a python script that read these values from the Watlow text
file temperature records. It was used in early stages of data analysis
to identify when Watlow errors within Labview had held the sample for
longer than the programmed step duration. When these discrepancies were
identified, the ``Nominal Time'' was updated to reflect the time the
sample was held at the programmed temperature. These times are not
directly used for any of the data reduction.

The ``Total Time'' column is the difference between the analysis
timestamp and the previous analysis timestamp. These timestamps are
recorded during the inlet of evolved, purified Helium into the
Helix-MC-Plus\textsuperscript{10K}, which immediately precedes the pump
out of the lamp cell chamber to the turbo pump. Therefore, the ``Total
Time'' represents the total amount of time that helium was allowed to
permeate through the silica glass ampoule into the UHV lamp cell,
including both the time the sample spent ``at temperature'' and the time
that the sample was cooling to or at room temperature (24-26°C). Since
the pump out times were only 1 minute and always occurred when the
sample was coolest (and permeation was slowest), permeation from the
ampoule during pump out is neglected in the data reduction. Less than
one thousandth of the helium remaining in the ampoule would permeate out
within one minute at room temperature (Eq.~\eqref{eq:naive}, Chapter~\ref{ch:apparatus}, with log
permeability (ln(K)) at 26°C of 23.88 ln(atoms/cm*atm*s)). Due to the
formatting of the time stamps, these values have a 1-minute resolution.

The ``Cooling Time'' is the difference between the ``Total Time'' and
``Nominal Time''. It is greater than or equal to the time to purify
evolved gas in preparation for analysis: \textasciitilde26 minutes. This
value is used for permeation corrections when low-temperature permeation
steps with cooling times greater than 26 minutes are used.

\subsection*{Abundances, 4He Blank Corrections, Measured
4He/3He}\label{abundances-4he-blank-corrections-measured-4he3he}

\subsubsection*{Sensitivity and Uncertainty
Calculations}\label{sensitivity-and-uncertainty-calculations}

3He and 4He sensitivities were calculated for each analysis of the MM
standard gas, assuming that each standard aliquot contained
157.50±0.41(2$\sigma$) nccSTP 4He and 3.567±0.040(2$\sigma$) pccSTP 3He. These values
were obtained from sensitivities from a known amount of a standard gas
with similar helium abundances and with the signals from 10 aliquots of
the new MM standard tank in June 2023.

1 sigma uncertainties related to internal reproducibility were estimated
for each sensitivity calculation as the rolling standard deviation of
sensitivities calculated from 5 MM standards. In other words, the
internal reproducibility of the sensitivity was estimated as the
standard deviation of each calculated sensitivity and the 2
sensitivities immediately before and after. Uncertainties related to
external reproducibility from the amount of 3He and 4He in each standard
aliquot were added to the total uncertainty as well.

Sensitivities and their uncertainties were interpolated to times between
standard analyses using a Piecewise Cubic Hermite Interpolating
Polynomial (PCHIP) from the SciPy package in python. These interpolated
sensitivities and uncertainties were used to convert the 3He and 4He
signals (cps and fA respectively) to absolute atomic abundances and
uncertainties in units of pccSTP and nccSTP respectively. The
sensitivities and interpolations are graphically displayed in the
``STDs'' tab in a plot generated using the Matplotlib package in python.

\subsubsection*{Helium-4 Blank Corrections}\label{helium-4-blank-corrections}

Throughout the experiment, low temperature permeation steps yielded
unpredictable 4He signals that did not exponentially decay like the 3He
signals did. This effect is most obvious when considering analyses where
cooling times were constant (\textasciitilde1560 seconds). In these
cases, 3He signals always decreased exponentially whereas the 4He
signals often increased (or decreased) erratically. A blank was measured
at the beginning of the experiment and showed negligible blank for 3He
(regressing to below 0 cps at the inlet time) and a very low value for
4He (4.84x10\textsuperscript{-2} fA). Low temperature permeation steps
sometimes yielded much higher signals despite failing to conform to the
expectation of exponential decay that is observed in the same steps' 3He
measurements. For instance, the 4He signal in a cool permeation step at
105°C from reference file He30509.txt was 7.93x10\textsuperscript{0} fA,
even when the high temperature experimental step preceding it
(He30502.txt) yielded a similar value: 8.28x10\textsuperscript{0} fA.
For the same pair of analyses, the 3He signal matched expectations and
varied by two orders of magnitude. 4He from permeation steps never fit
the equation for cumulative permeation, \(m_{cum} = n_{0}\left( 1 - e^{- \lambda t} \right)\)
(Eq.~\eqref{eq:cumyield}, Chapter~\ref{ch:apparatus}),
and thus could not be used to infer the permeability for 4He or for the
residual amount of 4He remaining in the ampoule (\(n_{0}\)). For
comparison, the same equation always fit the low temperature permeation
steps well for 3He when the cooling time of each analysis was constant.

The erratic behavior in the 4He data is best explained by a shifting
blank during the experiment, which was not measured or quantified
directly. It was largest in the middle of the experiment and returned to
low values by the end of the experiment that were similar to the blank
measured before the experiment. Because the blank for 4He was not well
controlled and appeared to have affected many of the values observed
even in the high temperature experimental steps, a blank correction was
deemed necessary. The 4He data of the permeation steps are interpreted
here to be dominated by blanks, so high temperature experimental steps
were coarsely blank corrected by subtracting the average 4He signal of
the subsequent permeation steps. The entirety of this correction was
conservatively added to the 2$\sigma$ uncertainty of the experimental data for
4He. One blank corrected measurement was within error of zero after
permeation corrections and was not plotted. The permeation lag was
estimated for the 4He data using the 3He data and assuming that the 4He
permeability was 0.91357±0.14603 (2$\sigma$) that of the 3He permeability.
These relative permeability values for 3He and 4He were measured at
Caltech for a similar silica glass ampoule half-filled with water and
equilibrated with \textasciitilde{} 0.1 torr of tank helium. This
process is described in greater detail in the Permeation Lag Corrections
section.

Some 3He permeation steps that followed especially low-yield
experimental steps were within error of baseline 3He cps values
(similarly to the measured blank at the beginning of the experiment).
These low measurements always preceded by measurements where 3He was
quantified above blank, so permeation lag corrections were calculated
using those values instead. In these cases, uncertainty from the
permeation lag corrections dominate the uncertainty of the
diffusivities, but do not affect the interpretation of the dataset.

\subsubsection*{Bracketing Method Ratios}\label{ratios}

4He/3He ratios are presented with reference to the irradiated Winsor
Goethite bulk ratio (\textasciitilde1000 or \textasciitilde1000 Ra) and
as standalone values. Uncertainties are propagated in quadrature from
the measured 3He and 4He abundances and from the bulk ratio, when
applicable. The blank and permeation corrected ratios have the same
general shape as the uncorrected ratios. The ratios plateau around a
cumulative 3He fraction of 0.06 at a bulk normalized 4He/3He ratio of
0.75. The early plateau in 4He/3He ratios with 3He fraction occurs
similarly in the vacuum step heat experiments of the same sample, in
agreement with previous studies \citep{monteiro2017,farley2024},
suggesting a relatively low loss of radiogenic helium for the Winsor
Goethite. The vacuum and hydrothermal experiments differ in the
bulk-normalized 4He/3He ratio at the plateau, with the hydrothermal
experiment plateauing at a lower 4He/3He ratio of 0.75 times the bulk
ratio, whereas the vacuum experiments plateau at the bulk ratio. This
could be explained by a 4He-retentive inclusion in the hydrothermal
experiment of this work similar to those discussed by \citet{vasconcelos2013}, or may reflect trapping of 3He in newly formed hematite
during the early steps of the vacuum experiment.

\subsection*{Permeation Lag Corrections for
3He}\label{permeation-lag-corrections-for-3he}

Helium evolved from the irradiated goethite must first permeate through
the silica glass ampoule into the high vacuum of the lamp cell to be
analyzed, a process which is non-instantaneous. Measurements can
therefore underestimate the helium production of the sample during a
step, as some of the produced helium will remain within the ampoule
after a high temperature step, and measurements can also overestimate
the amount of produced helium if there is some amount of helium
remaining from previous step heats within the ampoule at the beginning
of a new experimental step. To address this effect, we quantify the
amount of helium remaining within the ampoule between high temperature
experimental steps, and apply corrections using the rearranged form of
Eq.~\eqref{eq:permlag} (Chapter~\ref{ch:apparatus}), solved for \(m\):

\[\begin{array}{r}
m = n_{0} + \omega - n_{t}
\end{array}\]

Where \(\omega\) is the true helium production of the sample in a high
temperature experimental step, \(n_{0}\) is the amount of helium within
the ampoule at the beginning of the step, \(n_{t}\) is the amount of
helium within the ampoule at the end of the experiment, and \emph{m} is
the amount measured. Values of \(\omega\) and \emph{m} are most often
within 10\% of one another but differ by up to 47.2\% for 3He within
this dataset. \(n_{0}\) and \(n_{t}\) are tracked in the ``3He \(n_{0}\)
(pccSTP)'' column, where \(n_{t}\) for any measurement is the \(n_{0}\)
value of the next row.

At low temperatures (\textless120°C), when no new helium is evolved from
the sample, the measured value is simply the difference between the
helium contents of the ampoule before the measurement and after
(Eq.~\eqref{eq:permlag} with \(\omega = 0\)):

\[\begin{array}{r}
m = n_{0} - n_{t}
\end{array}\]

In this case, all the helium measured is a result of helium permeation
through the silica glass ampoule. Helium permeation is a first order
process that follows Eq.~\eqref{eq:naive} (Chapter~\ref{ch:apparatus}),
with \(\lambda = KGRT\):

\[\begin{array}{r}
n_{t} = n_{0}e^{- KGRTt}
\end{array}\]

Where \emph{K} is the temperature-dependent permeability coefficient,
\emph{G} is a geometric factor which can be determined through
measurements of the ampoule's inner and outer diameter, \emph{R} is the
gas constant, \emph{T} is the temperature in Kelvin, and \emph{t} is the
time in seconds. For a cylindrical ampoule, the value of G can be
calculated with measurements of the inner and outer radii using
Eq.~\eqref{eq:geomfactor} (Chapter~\ref{ch:apparatus}).
Inner and outer radii were measured using the amscope program at 11x
zoom setting. These radii were both measured on opposing sides of the
ampoule where there were air bubbles, allowing for the accurate sighting
of the inner walls of ampoule. The inner diameter, \emph{(2a)} was
measured at 505 and 540 microns; the outer diameter \emph{(2b)} was
measured at 787 and 832 microns (Fig.~\ref{fig:ch2-ampoule-dims}). The geometric factor \emph{G}
was calculated for both pairs of measurements and averaged to give a
value of \emph{G} of 6708.6 cm\textsuperscript{-1}.


For permeation steps, Eqs.~\eqref{eq:permlag} (with \(\omega=0\)) and
\eqref{eq:naive} apply because the temperature is too low to evolve any
helium from the sample, and the \(\omega=0\) form of
Eq.~\eqref{eq:permlag} can be substituted into Eq.~\eqref{eq:naive} to
relate a measurement to \(n_{0}\). This equation is calculated in the
column entitled ``3He n0 calc (pccSTP) (m/(1-e\^{}-$\lambda$t))''. Errors in the
measured value \emph{m} and the permeability \emph{K} are propagated in
quadrature. This is the informed-bracketing relation of
Eq.~\eqref{eq:n0informed} (Chapter~\ref{ch:apparatus}), expressed with
Eq.~\eqref{eq:permsol}'s exponential form for \(^{He}F\):

\[\begin{array}{r}
n_{0} = \frac{m}{\left( 1 - e^{- KGRT} \right)}
\end{array}\]

The ``3He \(n_{0}\) (pccSTP)'' column tracks the accepted value of
n\textsubscript{0} (and n\textsubscript{t}, which is n\textsubscript{0}
in the subsequent column) for ultimate use in permeation corrections
using Eq.~\eqref{eq:n0informed}. Depending on the step, this value is
determined one of three ways.

Most permeation steps in this study are ``Informed Bracket''
measurements, which set n\textsubscript{0} directly from their own
measurement using Eq.~\eqref{eq:n0informed} -- calculated in the
``3He n0 calc (pccSTP) (m/(1-e\^{}-$\lambda$t))'' column.

Permeation steps between Informed Bracket measurements instead track the
decreasing value of n\textsubscript{0} by the default first-order decay
of Eq.~\eqref{eq:naive}: the calculated n\textsubscript{0} value for step
\emph{i} is \(n_{0}^{i} = n_{0}^{i - 1}e^{- RG(KTt)_{i - 1}}\). This
default treatment is used whenever the previous step's cooling time
matched the \textasciitilde26-minute window used to calibrate the
effective permeabilities themselves (\S\ref{experimental-methods}), and
is generally preferred over re-inverting from an Informed Bracket
measurement because it lets n\textsubscript{0} be set by earlier
measurements with larger signals, avoiding error inflation from
near-blank permeation steps. Errors associated with \emph{K} and
\(n_{0}^{i - 1}\) are propagated through Eq.~\eqref{eq:naive} in
quadrature.

A small number of near-blank permeation steps, whose own signal was too
close to baseline to invert reliably, instead had n\textsubscript{0} set
by direct mass balance from an adjacent, reliably determined step:
forward, by subtracting that step's own measured amount \emph{m} from
the following step's n\textsubscript{0}; or backward, by adding it to
the preceding step's n\textsubscript{0}.

Finally, some permeation steps and series -- including the arbitrarily
long ``complete purge'' steps at the end of the experiment discussed in
\S\ref{complete-degas-test} -- were long enough to be treated as
emptying the ampoule entirely, so that n\textsubscript{0} for the
following row was set to 0 rather than inverted from
Eq.~\eqref{eq:naive}.

One further complication affects two adjacent high-temperature steps
(195°C/1106 min and 185°C/740 min): the Informed Bracket measurement
bracketing them cooled to room temperature and continued permeating for
far longer (\textasciitilde13 hours) than the \textasciitilde26-minute
window used to calibrate every other step's effective permeability, and
so would have needed a one-off, separately calibrated room-temperature
permeability not used anywhere else in this reduction. Rather than build
a special-case correction around two points, these two steps are instead
excluded from the fitted Stage 1 regression (\S\ref{ch2-si-distilled-naive}),
its Arrhenius plot, and the diffusivity comparison of
\S\ref{complete-degas-test}, and marked (not derived) with ``$^\dagger$''
in Table~\ref{tab:ch2-distilled-naive-3he} below: their differing cooling
time, rather than their signal strength, is the reason for the exclusion.


\section*{Na\"ive-Method Isothermal Permeability Dataset}\label{naive-apparent-lambdas-permeability-coefficient-description-and-methods}
\markright{\MakeUppercase{Na\"ive-Method Isothermal Permeability Dataset}}

The effective permeability coefficients used by the informed bracketing method below -- and the isothermal-permeation-fit methodology, cooling-tube numerical model, and Arrhenius regression used to determine them -- are described in full in the main text (\S\ref{experimental-methods}). The complete fitted dataset underlying that regression is reported here: Table~\ref{tab:ch2-naive-isothermal-raw} gives the raw per-step data for all 10 isothermal sequences, and Table~\ref{tab:ch2-naive-isothermal-fits} gives the lambda/$n_0$ fit to each sequence (Fig.~\ref{fig:ch2-naive-cumulative-yield}) that together define the Arrhenius regression of Fig.~\ref{fig:ch2-naive-arrhenius}.

\ChTwoNaiveRawTable

\clearpage

\ChTwoNaiveFitTable

\subsection*{Na\"ive-Method Step-Heat Dataset}\label{ch2-si-distilled-naive}

The complete corrected step-heat dataset underlying the na\"ive bracketing
method's apparent permeabilities (Fig.~\ref{fig:ch2-naive-arrhenius}) is reported
in Table~\ref{tab:ch2-distilled-naive-3he}, reporting the central value and
2$\sigma$ uncertainty for each quantity, analogous to
Table~\ref{tab:ch2-distilled-3he} in the
main text for the informed bracketing (``treatment as uncertainty'')
method. 4He is not treated by the na\"ive/informed bracketing method: 4He
permeation-step signals were blank level (\nameref{helium-4-blank-corrections}),
so no direct, step-by-step quantification of the ampoule's 4He inventory was
possible, and 4He throughout this chapter is instead reduced exclusively by
the bounding method of the main text (\S\ref{analytical-methods}). The
na\"ive/informed bracketing dataset reported here is retained only for 3He,
as the internal check on the bounding method discussed in
\S\ref{complete-degas-test}. ``--'' denotes steps for which a quantity could not
be determined (see \S\ref{finite-difference-diffusion-model-of-cooling-tube}
in the main text).

\begin{landscape}
\begin{footnotesize}
\begin{longtable}{r r r r r r r r r r r r}
\caption{$^3$He step-heat dataset for HT-Win-06-01b under the na\"ive bracketing method: measured and cumulative $^3$He, the initial and residual ampoule content ($n_{0}$, $n_{t}$), the true step production ($\omega$), and the apparent diffusivity $\ln(D/a^2)$, each with its 2$\sigma$ uncertainty. Six rows have derived values marked with ``$^*$'': these steps have measured $^3$He $<0.1$ pccSTP and are excluded as low-signal, as in Table~\ref{tab:ch2-distilled-3he}. Two of those six (195\textdegree C/1106 min and 185\textdegree C/740 min) are additionally marked ``$^\dagger$'': their $n_0$ depended on an atypically long room-temperature permeation rather than the standard cooling window used to calibrate the effective permeability (\S\ref{sec:ch2-si-detailed-methods}), so they are excluded (independently of the low-signal exclusion) from the fitted Stage 1 regression, Fig.~\ref{fig:ch2-bracket-arrhenius-zoom}, and Fig.~\ref{fig:ch2-diffusivity-comparison}.}\label{tab:ch2-distilled-naive-3he}\\
\toprule
Temp & Time & $^3$He meas. & Cum. $^3$He & $n_{0}$ & $\pm2\sigma$ & $n_{t}$ & $\pm2\sigma$ & $\omega$ & $\pm2\sigma$ & $\ln(D/a^2)$ & $\pm2\sigma$ \\
(°C) & (min) & (pccSTP) & (pccSTP) & (pccSTP) & & (pccSTP) & & (pccSTP) & & & \\
\midrule
\endfirsthead
\multicolumn{12}{l}{\textit{(continued)}} \\
\toprule
Temp & Time & $^3$He meas. & Cum. $^3$He & $n_{0}$ & $\pm2\sigma$ & $n_{t}$ & $\pm2\sigma$ & $\omega$ & $\pm2\sigma$ & $\ln(D/a^2)$ & $\pm2\sigma$ \\
(°C) & (min) & (pccSTP) & (pccSTP) & (pccSTP) & & (pccSTP) & & (pccSTP) & & & \\
\midrule
\endhead
\midrule
\multicolumn{12}{r}{\textit{continued on next page}} \\
\endfoot
\bottomrule
\endlastfoot
170 & 300 & 0.9630 & 0.0000 & \(0.00\!\times\!10^{0}\) & \(0.00\!\times\!10^{0}\) & 0.1968 & \(2.22\!\times\!10^{-2}\) & 1.1597 & \(2.63\!\times\!10^{-2}\) & -21.910 & 0.045 \\
180 & 463 & 0.5442 & 0.9630 & \(7.26\!\times\!10^{-2}\) & \(1.29\!\times\!10^{-2}\) & 0.0782 & \(4.93\!\times\!10^{-3}\) & 0.5497 & \(1.60\!\times\!10^{-2}\) & -22.184 & 0.089 \\
190 & 550 & 0.5292 & 1.6241 & \(1.95\!\times\!10^{-2}\) & \(2.54\!\times\!10^{-3}\) & 0.0679 & \(8.20\!\times\!10^{-3}\) & 0.5777 & \(1.30\!\times\!10^{-2}\) & -21.976 & 0.085 \\
200 & 500 & 0.5190 & 2.2112 & \(3.28\!\times\!10^{-2}\) & \(5.42\!\times\!10^{-3}\) & 0.0532 & \(5.77\!\times\!10^{-3}\) & 0.5394 & \(1.15\!\times\!10^{-2}\) & -21.702 & 0.088 \\
210 & 475 & 0.5770 & 2.7646 & \(1.41\!\times\!10^{-2}\) & \(2.11\!\times\!10^{-3}\) & 0.0552 & \(5.71\!\times\!10^{-3}\) & 0.6182 & \(1.17\!\times\!10^{-2}\) & -21.311 & 0.082 \\
220 & 450 & 0.6333 & 3.3799 & \(1.10\!\times\!10^{-2}\) & \(1.80\!\times\!10^{-3}\) & 0.0469 & \(4.58\!\times\!10^{-3}\) & 0.6692 & \(1.24\!\times\!10^{-2}\) & -20.991 & 0.080 \\
230 & 425 & 0.7899 & 4.0551 & \(5.87\!\times\!10^{-3}\) & \(1.86\!\times\!10^{-3}\) & 0.0639 & \(3.48\!\times\!10^{-3}\) & 0.8479 & \(1.25\!\times\!10^{-2}\) & -20.514 & 0.066 \\
225 & 908 & 0.7513 & 4.8885 & \(1.95\!\times\!10^{-2}\) & \(2.45\!\times\!10^{-3}\) & 0.0420 & \(4.94\!\times\!10^{-3}\) & 0.7738 & \(1.42\!\times\!10^{-2}\) & -21.200 & 0.121 \\
215 & 450 & 0.1415 & 5.6842 & \(1.31\!\times\!10^{-2}\) & \(1.85\!\times\!10^{-3}\) & 0.0278 & \(3.54\!\times\!10^{-3}\) & 0.1563 & \(6.48\!\times\!10^{-3}\) & -22.015 & 0.316 \\
205 & 520 & 0.0763 & 5.8529 & \(7.18\!\times\!10^{-3}\)$^*$ & \(1.23\!\times\!10^{-3}\)$^*$ & 0.0125$^*$ & \(1.63\!\times\!10^{-3}\)$^*$ & 0.0816$^*$ & \(8.10\!\times\!10^{-3}\)$^*$ & -22.788$^*$ & 0.829$^*$ \\
195 & 1106 & 0.0670 & 5.9491 & \(2.12\!\times\!10^{-3}\)$^{*\dagger}$ & \(4.02\!\times\!10^{-4}\)$^{*\dagger}$ & 0.0307$^{*\dagger}$ & \(2.86\!\times\!10^{-3}\)$^{*\dagger}$ & 0.0956$^{*\dagger}$ & \(3.23\!\times\!10^{-3}\)$^{*\dagger}$ & -23.370$^{*\dagger}$ & 1.094$^{*\dagger}$ \\
185 & 740 & 0.0227 & 6.0246 & \(1.25\!\times\!10^{-2}\)$^{*\dagger}$ & \(1.64\!\times\!10^{-3}\)$^{*\dagger}$ & 0.0090$^{*\dagger}$ & \(1.06\!\times\!10^{-3}\)$^{*\dagger}$ & 0.0192$^{*\dagger}$ & \(2.00\!\times\!10^{-3}\)$^{*\dagger}$ & -24.567$^{*\dagger}$ & 2.822$^{*\dagger}$ \\
175 & 920 & 0.0180 & 6.0610 & \(3.68\!\times\!10^{-3}\)$^*$ & \(5.02\!\times\!10^{-4}\)$^*$ & 0.0045$^*$ & \(4.42\!\times\!10^{-4}\)$^*$ & 0.0189$^*$ & \(7.42\!\times\!10^{-4}\)$^*$ & -24.796$^*$ & 3.797$^*$ \\
200 & 260 & 0.0224 & 6.0814 & \(6.53\!\times\!10^{-4}\)$^*$ & \(1.17\!\times\!10^{-4}\)$^*$ & 0.0095$^*$ & \(1.23\!\times\!10^{-3}\)$^*$ & 0.0313$^*$ & \(1.65\!\times\!10^{-3}\)$^*$ & -23.022$^*$ & 1.495$^*$ \\
210 & 520 & 0.0812 & 6.1081 & \(7.48\!\times\!10^{-3}\)$^*$ & \(9.98\!\times\!10^{-4}\)$^*$ & 0.0145$^*$ & \(1.32\!\times\!10^{-3}\)$^*$ & 0.0882$^*$ & \(3.85\!\times\!10^{-3}\)$^*$ & -22.669$^*$ & 1.135$^*$ \\
220 & 1167 & 0.5766 & 6.1914 & \(2.82\!\times\!10^{-3}\) & \(1.44\!\times\!10^{-3}\) & 0.0276 & \(3.24\!\times\!10^{-3}\) & 0.6014 & \(2.45\!\times\!10^{-2}\) & -21.504 & 0.184 \\
230 & 370 & 0.3577 & 6.7796 & \(8.68\!\times\!10^{-3}\) & \(3.96\!\times\!10^{-3}\) & 0.0422 & \(5.13\!\times\!10^{-3}\) & 0.3912 & \(1.48\!\times\!10^{-2}\) & -20.712 & 0.146 \\
240 & 320 & 0.6767 & 7.1563 & \(1.72\!\times\!10^{-2}\) & \(2.42\!\times\!10^{-3}\) & 0.0826 & \(9.94\!\times\!10^{-3}\) & 0.7420 & \(1.57\!\times\!10^{-2}\) & -19.807 & 0.099 \\
245 & 250 & 0.5549 & 7.8582 & \(2.13\!\times\!10^{-2}\) & \(8.87\!\times\!10^{-3}\) & 0.0516 & \(5.65\!\times\!10^{-3}\) & 0.5852 & \(1.46\!\times\!10^{-2}\) & -19.710 & 0.125 \\
250 & 210 & 0.6934 & 8.4743 & \(1.08\!\times\!10^{-2}\) & \(4.52\!\times\!10^{-3}\) & 0.0718 & \(8.15\!\times\!10^{-3}\) & 0.7545 & \(1.41\!\times\!10^{-2}\) & -19.200 & 0.102 \\
245 & 270 & 0.4659 & 9.2086 & \(2.48\!\times\!10^{-2}\) & \(3.48\!\times\!10^{-3}\) & 0.0498 & \(5.84\!\times\!10^{-3}\) & 0.4909 & \(1.01\!\times\!10^{-2}\) & -19.810 & 0.197 \\
240 & 320 & 0.2857 & 9.7206 & \(2.07\!\times\!10^{-2}\) & \(2.81\!\times\!10^{-3}\) & 0.0339 & \(4.08\!\times\!10^{-3}\) & 0.2989 & \(6.58\!\times\!10^{-3}\) & -20.433 & 0.312 \\
235 & 370 & 0.2199 & 10.0357 & \(1.66\!\times\!10^{-2}\) & \(2.20\!\times\!10^{-3}\) & 0.0273 & \(3.29\!\times\!10^{-3}\) & 0.2307 & \(5.23\!\times\!10^{-3}\) & -20.810 & 0.401 \\
230 & 420 & 0.1400 & 10.2736 & \(1.32\!\times\!10^{-2}\) & \(1.76\!\times\!10^{-3}\) & 0.0255 & \(2.53\!\times\!10^{-3}\) & 0.1523 & \(4.00\!\times\!10^{-3}\) & -21.332 & 0.603 \\
225 & 652 & 0.1494 & 10.4279 & \(1.26\!\times\!10^{-2}\) & \(1.62\!\times\!10^{-3}\) & 0.0131 & \(1.47\!\times\!10^{-3}\) & 0.1499 & \(3.42\!\times\!10^{-3}\) & -21.772 & 0.747 \\
230 & 500 & 0.1550 & 10.5908 & \(7.34\!\times\!10^{-4}\) & \(1.67\!\times\!10^{-3}\) & 0.0128 & \(1.32\!\times\!10^{-3}\) & 0.1671 & \(3.37\!\times\!10^{-3}\) & -21.383 & 0.461 \\
235 & 450 & 0.1951 & 10.7582 & \(2.57\!\times\!10^{-3}\) & \(4.18\!\times\!10^{-4}\) & 0.0334 & \(1.85\!\times\!10^{-3}\) & 0.2259 & \(3.29\!\times\!10^{-3}\) & -20.956 & 0.364 \\
240 & 400 & 0.2425 & 10.9638 & \(6.90\!\times\!10^{-3}\) & \(1.45\!\times\!10^{-3}\) & 0.0251 & \(3.16\!\times\!10^{-3}\) & 0.2607 & \(1.23\!\times\!10^{-2}\) & -20.672 & 0.318 \\
245 & 350 & 0.3150 & 11.2328 & \(5.79\!\times\!10^{-3}\) & \(1.05\!\times\!10^{-3}\) & 0.0395 & \(4.64\!\times\!10^{-3}\) & 0.3487 & \(9.39\!\times\!10^{-3}\) & -20.219 & 0.240 \\
250 & 300 & 0.3716 & 11.5689 & \(1.65\!\times\!10^{-2}\) & \(2.24\!\times\!10^{-3}\) & 0.0467 & \(5.24\!\times\!10^{-3}\) & 0.4018 & \(9.62\!\times\!10^{-3}\) & -19.888 & 0.209 \\
250 & 300 & 0.3640 & 11.9627 & \(2.43\!\times\!10^{-2}\) & \(4.37\!\times\!10^{-3}\) & 0.0339 & \(3.82\!\times\!10^{-3}\) & 0.3735 & \(7.99\!\times\!10^{-3}\) & -19.927 & 0.242 \\
250 & 300 & 0.3452 & 12.3490 & \(1.13\!\times\!10^{-2}\) & \(1.60\!\times\!10^{-3}\) & 0.0395 & \(4.48\!\times\!10^{-3}\) & 0.3733 & \(7.91\!\times\!10^{-3}\) & -19.895 & 0.244 \\
250 & 10250 & 8.4829 & 12.7167 & \(1.91\!\times\!10^{-2}\) & \(2.53\!\times\!10^{-3}\) & 0.1165 & \(7.38\!\times\!10^{-3}\) & 8.5804 & \(1.28\!\times\!10^{-1}\) & -19.958 & 0.189 \\
240 & 785 & 0.2261 & 21.2187 & \(2.90\!\times\!10^{-9}\) & \(6.85\!\times\!10^{-9}\) & 0.0263 & \(2.92\!\times\!10^{-3}\) & 0.2525 & \(2.21\!\times\!10^{-2}\) & -20.660 & 0.580 \\
\midrule
\textbf{Total} &  &  & 146.1483 &  &  &  &  &  &  &  &  \\
\end{longtable}
\end{footnotesize}
\end{landscape}

% 4He na\"ive-bracketing step-heat table (tab:ch2-distilled-naive-4he) and the
% derived 4He/3He ratio table (tab:ch2-distilled-naive-ratios) removed here:
% both depended on the "Permeation Lag Corrections for 4He" reduction removed
% above. 4He is reduced exclusively by the bounding method in the main text.

Beyond the ampoule's residual helium content itself (\S\ref{complete-degas-test}), the two methods' derived apparent diffusivities agree just as closely. Fig.~\ref{fig:ch2-diffusivity-comparison} compares $\ln(D/a^2)$ from the permeation-step method against the bounding method's central value for every step in the experiment, with the bounding method's own ``treatment as uncertainty'' spread on one axis and the permeation-step method's propagated uncertainty on the other; the two steps marked ``$^\dagger$'' in Table~\ref{tab:ch2-distilled-naive-3he} above are excluded for the reason given there. The two diffusivity datasets track the 1:1 line within uncertainty across the full range of the experiment. Fitting the permeation-step method's own diffusivities with the same Stage 1 procedure used for the bounding method in the main text (Fig.~\ref{fig:ch2-bracket-arrhenius-zoom}) yields an activation energy of \ChTwoBracketEa$\pm$\ChTwoBracketEaErr\ kJ/mol (R\textsuperscript{2}=\ChTwoBracketRtwo, n=\ChTwoBracketN) and a closure temperature of \ChTwoBracketTc$\pm$\ChTwoBracketTcErr\textdegree C, indistinguishable from the bounding method's own \ChTwoEa$\pm$\ChTwoEaErr\ kJ/mol and \ChTwoTc$\pm$\ChTwoTcErr\textdegree C (Fig.~\ref{fig:ch2-arrhenius-zoom}, main text).

\begin{figure}[htbp]
\centering
\includegraphics[width=0.8\textwidth]{Winsor_diffusivity_naive_vs_bounding.pdf}
\caption{Apparent \textsuperscript{3}He diffusivities for HT-Win-06-01b from the permeation-step method ($y$-axis) against the bounding method's central value ($x$-axis, Eq.~\eqref{eq:steadystate}), with the bounding method's ``treatment as uncertainty'' LB/UB spread and the permeation-step method's own propagated uncertainty plotted on their respective axes. The dashed line is 1:1. Two steps (195\textdegree C/1106 min and 185\textdegree C/740 min) are excluded for the reason given in Table~\ref{tab:ch2-distilled-naive-3he}'s caption.}
\label{fig:ch2-diffusivity-comparison}
\end{figure}

\begin{figure}[htbp]
\centering
\includegraphics[width=0.8\textwidth]{winsor_goethite_Diffusivities_Arrhenius_Zoomin_naive.pdf}
\caption{Analogue of Fig.~\ref{fig:ch2-arrhenius-zoom} (main text) using the permeation-step method's own \textsuperscript{3}He diffusivities and its own Stage 1 linear regression, rather than the bounding method's. The resulting activation energy and closure temperature (\ChTwoBracketEa$\pm$\ChTwoBracketEaErr\ kJ/mol; \ChTwoBracketTc$\pm$\ChTwoBracketTcErr\textdegree C) agree closely with the bounding method's own (\ChTwoEa$\pm$\ChTwoEaErr\ kJ/mol; \ChTwoTc$\pm$\ChTwoTcErr\textdegree C).}
\label{fig:ch2-bracket-arrhenius-zoom}
\end{figure}

\clearpage

The finite-difference numerical model used to evaluate the contribution of the 26-minute cooling period to these apparent permeabilities -- and to confirm that the ampoule's markedly elevated permeability is not an artifact of that cooling history -- is described in full in the main text (\S\ref{finite-difference-diffusion-model-of-cooling-tube}).

\section*{\texorpdfstring{Diffusivities
}{Diffusivities }}\label{diffusivities}
\markright{\MakeUppercase{Diffusivities}}

Apparent values of D/a\textsuperscript{2} for 3He and 4He were
calculated from times and corrected cumulative yields using the
equations of \citet{fechtigkalbitzer1966}. Uncertainties in the corrected
values are propagated into the cumulative yield, and all errors are
treated as uncorrelated for simplicity. Error bars plotted in the
Arrhenius plots on this tab are likely exaggerated. Derivatives
necessary to calculate uncertainty in the diffusivities using quadrature
are tabulated to the right of the sheet. The Stage 1 reference-line linear regression (Fig.~\ref{fig:ch2-arrhenius-zoom}),
its statistics, and the treatment of total-yield uncertainty in the regression are described in the main text
(\S\ref{analytical-methods}); the same regression underlies both the main text's treatment-as-uncertainty results
and the naive-bracketing dataset reported here.

\section*{Totals}\label{totals}
\markright{\MakeUppercase{Totals}}

The main text (\S\ref{experimental-methods}) already summarizes this final furnace-fusion step and its blank/re-extract
corrections; the additional spreadsheet-level detail below (volume-split calibration, standard-averaging convention)
is retained here for completeness. For the final step, the glass ampoule was broken and the remainder of
the helium was completely degassed in a resistance furnace at 1100°C.
The gas from this final step was split within the vacuum line on
calibrated volumes so that purified helium from less than 1\% of the
evolved gas was inlet into a Thermo Fisher Helix SFT mass spectrometer.
A gas standard with a known amount of 3He and 4He was measured twice
preceding the final fusion measurement of the hydrothermal Win-06-01b,
as was a split furnace hot blank. The gas standards were used to
calculate sensitivities, which were taken to be the mean of the two
measurements, with 1$\sigma$ uncertainty approximated by their standard
deviation. After the split measurement of the Win-06-01b ampoule, the
furnace was reheated to 1100°C to ensure complete extraction of the
helium in the sample. The final step was then taken to be the measured
value minus the hot blank, plus the re-extract, all divided by their
respective volume split fractions and multiplied by the calculated
sensitivity. Uncertainties were propagated in quadrature.


